\chapter{Related Work}
\label{ch:related-work}

This chapter reviews the automated C-to-Rust translation systems and the ways in which their translations are evaluated. It examines what kind of evidence each evaluation approach provides about a translation and where that evidence stops. It then considers studies that have independently reassessed published results, and from these two strands it derives the research gap that this thesis addresses.

\section{Evaluation Approaches in C-to-Rust Translation}
\label{sec:rw-approaches}\label{sec:rw-translation}

Automated C-to-Rust systems differ considerably in how they produce
translations. Rule-based tools such as C2Rust mechanically transform C
code into Rust, often retaining substantial use of \texttt{unsafe};
Laertes subsequently refines raw pointers into Rust references where
possible, using compiler feedback during the conversion process
\cite{c2rust,emre2021laertes}. LLM-assisted systems range from
approaches that refine C2Rust output, such as C2SaferRust, to
feedback-driven translation systems such as SmartC2Rust, which use
compilation failures, run-time errors, execution-flow discrepancies, and
assertion failures to guide revisions
\cite{nitin2025c2saferrust,shiraishi2026smartc2rust}. Other
configurations make a single model call for a project, whereas the ACTOR
study evaluates general-purpose coding agents, including Kiro, Claude
Code, and Codex, that can interact with development tools throughout
translation \cite{schesch2026actor}. These production mechanisms provide context; the present review focuses on the evidence used to judge their outputs.

Three components are useful for distinguishing evaluation procedures:
the \emph{corpus}, which defines the programs to be translated; the
\emph{oracle}, which determines whether a candidate meets the selected
acceptance criteria; and the \emph{metric}, which aggregates judgments
into a reported quantity such as compilation rate or test pass rate.
These components should not be conflated. Two studies may use similarly
named correctness metrics while selecting different programs, imposing
different acceptance conditions, or omitting different classes of
results from their denominators.

C-to-Rust benchmarks illustrate the importance of specifying the
evaluation population. CRUST-Bench comprises 100 C repositories with
Rust interfaces and functional tests \cite{khatry2025crustbench}. The
TRACTOR Public Test Corpus contains executable C programs and libraries
accompanied by test vectors for evaluating translations
\cite{tractor-test-corpus}. Sharing a benchmark name does not
necessarily mean sharing an identical experimental population: TRAIL and
DepTrans evaluate all 100 CRUST-Bench repositories, whereas ORBIT
evaluates 24 repositories, preferentially selecting larger programs; 22
of its 24 selections exceed 1,000 lines of code
\cite{wang2026trail,feng2026deptrans,farrukh2026orbit}. A reported
percentage is interpretable only in relation to the cases from which it
was calculated.

The surveyed methods draw on six overlapping categories of evidence,
summarized in Table~\ref{tab:rw-evidence}. Some of these categories describe a judgment
mechanism, whereas others describe how test inputs are obtained. In
particular, generating a test is not itself a behavioral oracle: the
generated test still requires an expected result or another comparison
procedure. A system may therefore use compilation checks, existing
tests, generated tests, and differential testing within a single
pipeline.

\begin{table}[htbp]
\centering\footnotesize
\caption{Evaluation approaches in C-to-Rust studies and limitations
documented in prior work. Rows follow the order of
Sections~\ref{sec:rw-compilation}--\ref{sec:rw-formal}. Citations after
each approach are examples of studies that use it rather than a census,
and one study may appear in several rows.}
\label{tab:rw-evidence}
\begin{tabularx}{\textwidth}{@{}>{\raggedright\arraybackslash}p{.27\textwidth}>{\raggedright\arraybackslash}X@{}}
\toprule
\textbf{Approach} & \textbf{Documented limitation} \\
\midrule
Compilation \cite{bedell2026llm4c2rust,feng2026deptrans,farrukh2026orbit} & With a compile-only reward, DepTrans reaches 97.9\% compilation and 0\% functional accuracy \cite{feng2026deptrans} \\
\addlinespace
Existing tests \cite{farrukh2026orbit,wang2026trail,feng2026deptrans,rustify2026,schesch2026actor} & PASS: Reboot's translations pass all supplied tests but 61.58--91.78\% of held-out tests \cite{wang2026reboot}. FAIL: 21 of 26 examined CRUST-Bench projects have expected values that cannot be derived from the C source \cite{schesch2026actor} \\
\addlinespace
Generated or ported tests \cite{schesch2026actor,leanh2026rustprint,yuan2026ptrtrans} & RustPrint's test translator tends to omit tests for features its own translation lost \cite{leanh2026rustprint} \\
\addlinespace
Reference similarity and model judgment \cite{feng2026deptrans,rustify2026,wang2026evoc2rust,leanh2026rustprint} & No program is executed; EvoC2Rust compares generated lines with a reference translation \cite{wang2026evoc2rust} \\
\addlinespace
Differential testing against C \cite{eniser2024flourine,sim2026encrust,shiraishi2026smartc2rust,actor2025} & Evidence is limited to executed inputs; Flourine accepts after five minutes without a counterexample \cite{eniser2024flourine}; C-side UB leaves no reference \\
\addlinespace
Symbolic and formal methods \cite{yang2024vert,bai2025rustassure} & VERT's full verification does not scale to loops bounded by input; RustAssure compares symbolic expressions by graph edit distance \cite{yang2024vert,bai2025rustassure} \\
\bottomrule
\end{tabularx}
\end{table}

The distinction between these sources matters for this thesis. The
reported decisions under study were produced by a build-and-test
procedure that compares translated programs with stored expected
results. The thesis subsequently examines those decisions using an
independent behavioral comparison against the C programs. The following
sections explain the evidence each procedure offers and the limits of
the conclusions drawn from it.

\section{Compilation and Existing Test-Suite Evaluation}
\label{sec:rw-compilation}\label{sec:rw-evaluation}

Compilation is a common condition in C-to-Rust evaluation, but a
successful build establishes only that the generated program meets the
requirements enforced by the compiler. Some studies evaluate
translations primarily at this level. LLM4C2Rust, for example, examines
\texttt{rustc} diagnostics and counts \texttt{unsafe} constructs without
running the translated programs as functional tests
\cite{bedell2026llm4c2rust}. The distinction between compilation and
correct behavior is especially visible in a DepTrans ablation: when its
reward was based only on compilation, the reported compilation success
rate reached 97.9\%, while functional accuracy was 0\%, because the
model could produce trivial function bodies invoking \texttt{panic!}
\cite{feng2026deptrans}. Compilability is consequently a useful
prerequisite for execution, but not evidence that a translation
preserves the behavior of its C source.

Many studies add functional tests to the build requirement. ORBIT uses
\texttt{cargo\ check} and \texttt{cargo\ test} against the interfaces
and tests supplied by CRUST-Bench \cite{farrukh2026orbit}. A comparison
table in RustAssure identifies seven of thirteen surveyed C-to-Rust
systems as using existing concrete test cases, whereas RustAssure itself
is classified under differential symbolic testing
\cite{bai2025rustassure}. In the TRACTOR Public Test Corpus, test
vectors contain inputs and expected outputs, and the corpus checks in its
continuous integration that every vector matches the behavior of the C
program. The runner evaluates Rust against these recorded expectations,
rather than executing C and Rust side by side when judging each
translation \cite{tractor-test-corpus,okhravi2026tractorbenchmark}. The
detailed vector semantics and runner behavior needed to reproduce the
original decisions are described in Chapter~\ref{ch:methodology}. At this level, a PASS
establishes conformity to the particular build requirements and test
expectations applied; its strength as behavioral evidence depends on the
content and adequacy of those tests.

One limitation is that a translation can pass the available tests while
differing from the source on untested inputs. Reboot provides direct
evidence: all six translated interpreters pass the test suites supplied
during translation, but their pass rates range from 61.58\% to 91.78\%
on separate validation suites written by the researchers and withheld
from the translation process \cite{wang2026reboot}. Related evidence
concerns the amount of functionality that available tests actually
exercise. In ORBIT's LTRE ablation, a basic coding-agent configuration
successfully built and passed the tests it executed while implementing
only 64.4\% of the functions and running 3.0\% of the original tests
\cite{farrukh2026orbit}. This result shows why coverage and functional completeness should be considered alongside the pass status. At function level, RustAssure also reports that the
existing tests for the \texttt{libcsv} and \texttt{u8c} libraries found
only one of the translation bugs exposed by its symbolic comparison
\cite{bai2025rustassure}. Together these findings show that passing a
finite suite may leave substantial behavior unexamined.

The reverse problem arises when a translation fails a test whose
expectation is unsuitable as a behavioral reference. CRUST-Bench issue
37 documents inconsistent expectations for the \texttt{Holdem\_Odds}
project: some hands that the C implementation identifies as a straight
flush are tested against expectations for a different poker-hand
category \cite{crustbench-issues}. Schesch and Ernst similarly report
that, among 26 CRUST-Bench projects they examined, 21 had tests whose expected values could not be derived from the C source. They separately exclude 13 of the benchmark's 100 projects as having unsatisfiable tests
\cite{schesch2026actor}. These observations are of different kinds: an expected value that cannot be derived from the C source is one problem, and a test that no translation can satisfy is another. Nevertheless, a FAIL by itself
establishes disagreement with the test criterion, not necessarily a
fault in the Rust translation relative to C.

Limitations of individual judgments also affect how aggregate results
are interpreted. A compilation rate may be computed per file, function,
module, or repository, and a behavioral score may be restricted to
translations that have already compiled
\cite{rustify2026,sim2026encrust,wang2026trail,wang2026evoc2rust}.
C2RustXW, for instance, computes its semantic-correctness measures only
for compilable Rust programs \cite{yan2026c2rustxw}. LAC2R evaluates
selected subsets of its benchmarks while retaining programs that fail to
compile in the denominator used for its reported rates
\cite{sim2025lac2r}. Test availability during translation introduces
another relevant condition: ACTOR reports 56 of 87 CRUST-Bench projects
passing when agents cannot see the benchmark tests and 82 of 87 when
they are allowed to repair against them \cite{schesch2026actor}. These
quantities are meaningful for their stated settings, but should not be
compared as if they represented identical populations and evaluation
procedures.

Thus, the build-and-test approach has three distinct limits for
interpreting published results. A PASS can overlook a behavioral
divergence that available tests do not exercise; a FAIL can arise from
an expectation that does not faithfully describe the C program; and a
reported success rate depends on how cases and decisions enter its
denominator. The first two issues motivate the two directions of
decision re-evaluation examined by RQ1 and RQ2, while the third matters
when the resulting evidence is used to reassess per-system success
rates.

\section{Automatically Generated and Translated Tests}
\label{sec:rw-generated-tests}

When a benchmark lacks tests suitable for generated Rust code, a
translation system may synthesize tests or convert tests from the
original C project. Such techniques can extend the set of executed
inputs, but generating inputs and establishing their expected outcomes
are separate tasks. More tests do not automatically imply stronger
evidence about behavioral preservation; the evidential value depends on
what the tests cover and how their results are judged.

ACTOR allows coding agents to construct tests and check candidate
translations against the C source \cite{schesch2026actor}. RustPrint
introduces a TestTranslator agent that ports maintainers' C tests to
Rust \cite{leanh2026rustprint}. PtrTrans combines LLMs with AFL++: a
model constructs a test harness, fuzzing collects C input-output pairs,
and selected pairs are converted into Rust unit tests
\cite{yuan2026ptrtrans}. These systems automate test construction in
different ways. In particular, obtaining expected behavior from an
actual C execution provides a different kind of evidence from having an
LLM infer what a program ought to return.

Test construction can also depend on the translation being evaluated.
RustPrint reports \emph{self-alignment bias}: a TestTranslator tends to
omit tests associated with features that its own translation failed to
preserve. The study addresses this concern through a cross-test protocol
that applies tests prepared by one system to translations produced by
others \cite{leanh2026rustprint}. This is an important methodological response to that concern. If a
system influences both the translation and the test cases used to
validate it, a positive result may reflect its ability to select
favorable cases as well as its ability to implement the required
functionality.

The independence of a generated suite must therefore be assessed in
terms of input provenance, expected-result construction, and the use of
testing feedback to revise or select translations. Automatically
generated tests can provide independent behavioral evidence when their
expected results are obtained from the C reference and the tested
translations are held fixed. Conversely, a test that was absent from the
benchmark may still be part of the translation process if its outcome is
used to drive repair. The present thesis adopts a post-publication
arrangement: its independent oracle does not inform the creation or
refinement of the released translations whose reported decisions it
re-examines.

\section{Reference-Similarity and Model-Based Judgments}
\label{sec:rw-model}

Other C-to-Rust evaluations use code-reference similarity or a language
model's judgment rather than executing the translated program. Although
both avoid the need to run every candidate, they should be
distinguished. A similarity metric scores features of the generated
representation against a reference; a model-based judge assesses a
supplied artifact against criteria expressed in a prompt or evaluation
procedure.

CodeBLEU extends token-level overlap measures with information about
syntax and data flow, and can be computed without executing the code
\cite{ren2020codebleu}. DepTrans and Rustify use CodeBLEU where
reference Rust implementations are available
\cite{feng2026deptrans,rustify2026}. EvoC2Rust reports an accepted-line
rate (AccRate) that compares generated lines with a reference translation produced with an LLM and then checked by people \cite{wang2026evoc2rust}. These measures can
be informative about structural resemblance or the amount of subsequent
editing. Nevertheless, two implementations can preserve the same
behavior while having very different syntax, and a nearly identical
condition or expression can still cause a meaningful behavioral
difference.

RustPrint uses CodeWikiBench and an LLM-as-judge procedure to compare
documentation generated from the C and Rust repositories. Its
functional-coverage measure rests on the assumption that similarity
between their documentation is informative about similarity of the
implemented functionality \cite{leanh2026rustprint}. DepTrans also uses
an LLM to assess C and Rust during translation refinement, but such
feedback inside a repair loop is not the same as a final independent
evaluation metric \cite{feng2026deptrans}. For either mechanism, the
conclusion remains conditional on the representational proxy or the
criteria available to the model. Neither directly establishes that the
two executable programs produce the same observation for a specified
input. The thesis therefore bases its own re-evaluation on comparisons of executed behavior.

\section{Differential Testing Against the C Program}
\label{sec:rw-differential}\label{sec:rw-equivalence}

Differential testing compares the observable outcomes of comparable
programs on common inputs \cite{mckeeman1998differential}. For C-to-Rust
translation, the original C program can serve as a behavioral reference
for the translated program without requiring an expected output to be
written for every new input. Inputs may come from a benchmark or be
generated by fuzzing. Coverage-guided fuzzers such as AFL++ mutate
candidate inputs and retain those that exercise previously unseen
program paths or transitions \cite{fioraldi2020aflpp,aflpp-docs}. If the
two executions differ according to the selected observation relation,
the input exposes a \emph{divergence}. Such a divergence is evidence of
disagreement, but it is not automatically evidence of a translation
error: the C reference execution and the comparison itself must also be
valid. Conversely, failing to find a divergence among finitely many
tests cannot establish equivalence on all possible inputs.

Several C-to-Rust systems combine differential testing with automated
input generation. Flourine uses cross-language differential fuzzing and
serialized program state across an FFI boundary to compare C and Rust;
its acceptance criterion treats the absence of a counterexample during a
five-minute fuzzing budget as practical equivalence, not as a
mathematical proof \cite{eniser2024flourine}. ENCRUST uses AFL++ to
generate 1,000 inputs for each of the Laertes libraries it evaluates. It
first executes the C implementation, discards inputs for which C crashes
or produces nondeterministic output across repeated runs, and uses the
remaining C observations to compare the Rust translation
\cite{sim2026encrust}. These approaches obtain behavioral reference
information directly from the program being translated, although their
conclusions remain limited by their inputs and observation models.

The observable behavior used in these comparisons also varies. RustMap
compares execution state at corresponding breakpoints to help localize
the point at which C and Rust diverge \cite{cai2025rustmap}. SmartC2Rust
records C outputs in advance to compare entry-point results without
requiring a C-to-Rust FFI call during testing. It additionally compares
execution-flow information, including call stacks and function inputs
and returns, to diagnose differences between implementations
\cite{shiraishi2026smartc2rust}. Final-output comparisons are relatively
straightforward for complete programs, whereas intermediate-state or
trace comparisons can help identify the location and nature of a
discrepancy. The latter require additional assumptions about which
states or events correspond across the two languages.

ACToR, developed by a group at the National University of Singapore,
pairs a translation agent with a discriminator that constructs
differential fuzzers to find inputs exposing mismatches between C and
Rust. Counterexamples are fed back into the translation loop for
refinement; the system also uses C-versus-Rust execution when assessing
its resulting translations \cite{actor2025}. Its equivalence objective
is formulated over a \emph{safe universe} of inputs considered valid for
the C program. ACToR should not be confused with the similarly named
ACTOR study by Schesch and Ernst on general-purpose coding agents and
TRACTOR. ACToR's additional evaluation of C2SaferRust translations is
discussed separately in Section~\ref{sec:rw-reevaluation} because it is a particularly
relevant precedent for assessing outputs from another translation
system.

\phantomsection\label{sec:bg-ub}%
A fundamental challenge for a C-based oracle is \emph{undefined
behavior} (UB). Under the C standard, undefined behavior imposes no
requirements on the outcome of the execution \cite{wg14n2176}. Signed
integer overflow, invalid memory access, and use of an object after its
lifetime ends are familiar examples. If the C program encounters UB on
an input, its observed output is not a language-defined reference to
which the Rust program must conform. A mismatch with that output is
therefore insufficient to show that Rust is incorrect, while a match is
not proof of correctness. This is different from inadequate test
coverage: even when both results are available, the supposed reference
may not have the required semantic status.

Existing evaluations handle this issue through different restrictions.
TRACTOR skips vectors explicitly labeled \texttt{has\_ub}
\cite{tractor-test-corpus}. ENCRUST discards inputs leading to C crashes
or nondeterministic results, but neither property coincides exactly with
undefined behavior \cite{sim2026encrust}. AddressSanitizer (ASan) and
UndefinedBehaviorSanitizer (UBSan) can detect selected memory errors and
undefined operations in instrumented C executions
\cite{asan-docs,ubsan}. A positive report of a covered UB condition
disqualifies that execution as a reliable reference; the absence of a
report, however, does not prove that no UB occurred. Sanitizer-based
screening consequently strengthens the interpretation of differential
results only within the classes of errors and executed operations that
the instrumentation can detect.

Differential testing thus provides empirical evidence within a defined
observation relation and execution budget. The thesis applies this principle after translations and their original decisions have been published, and it leaves the translations unchanged. Chapter~\ref{ch:methodology} defines its observation model, input generation,
stability checks, and rules for adjudicating divergent executions and
C-side UB.

\section{Symbolic and Formal Approaches}
\label{sec:rw-formal}

Translation validation provides a complementary approach: it checks
whether an individual translation is equivalent to its source, rather
than proving that every output of the translator must be correct
\cite{pnueli1998tv,necula2000tv}. Symbolic execution and formal
verification can reason about multiple program paths through
constraints, potentially providing stronger guarantees than a finite
collection of concrete tests. The conclusions they support nevertheless
depend on the semantics represented by their models, the assumptions
imposed on inputs and environments, and the bounds or proof obligations
actually discharged.

VERT generates Rust translations with an LLM and constructs a separate
Rust reference by lifting the original program through WebAssembly. It
combines property-based testing with bounded and, where feasible, full
verification using Kani \cite{yang2024vert}. The authors report
difficulty scaling full verification to programs containing loops whose
iteration counts depend on input. Consequently, a passed property-based
test, a bounded model-checking result, and the discharge of a full proof
obligation should not be interpreted as interchangeable evidence, even
when they arise within the same verification framework.

RustAssure takes a different symbolic approach at the function level. It
symbolically executes the C and Rust functions using KLEE, then compares
symbolic return-value expressions by graph edit distance
\cite{bai2025rustassure}. The comparison can reveal differences missed
by existing concrete tests, but similarity between selected symbolic
expressions does not automatically amount to a proof of behavioral
equivalence for a whole program. The unit of analysis, modeled outputs,
and approximation used in comparison all affect the scope of the result.

These methods extend the choices available for evaluating translations,
without eliminating the need to specify what is being checked. The
present thesis instead pursues a uniform empirical re-evaluation of a
defined population of previously released translations. This choice sets the scope of the evidence the thesis can report.

\section{Independent Re-evaluation of Reported Results}
\label{sec:rw-reevaluation}\label{sec:rw-threats}

The validity of test-based acceptance decisions has also been
investigated outside program translation. These studies provide
precedents for testing whether published evaluation outcomes are
supported by stronger or independently obtained evidence. They are relevant because they distinguish a system's originally reported success from the outcome of a
subsequent assessment. It is nevertheless necessary to distinguish
post-hoc audits of prior decisions from independent primary evaluations
and from tests performed to improve a system's own outputs. The studies
are reviewed in three groups: audits of acceptance decisions, evidence
about rejections and about C-to-Rust evaluations, and the closest
C-to-Rust precedents.

A first group of studies shows that published acceptance decisions can
be unreliable. In automated program repair, Qi et al.~re-examined
patches reported by
GenProg, RSRepair, and AE on a benchmark of 105 C defects. They found
that many patches did not pass even the validation tests used in the
original studies when their actual outputs were checked. Of 414 reported
GenProg patches, only 110 satisfied the test suite according to the
re-evaluation; subsequent analysis showed that truly correct repairs
were much rarer still \cite{qi2015plausibility}. This result exposes a
problem more fundamental than insufficient test coverage: the original
test infrastructure could accept patches despite their failure on
supplied test cases. Smith et al.~showed, with tests held out from
repair, that patches passing the repair-time suite need not behave
correctly on independent tests \cite{smith2015overfitting}. These studies question
positive decisions, but they do not systematically reconsider previously
rejected candidates.

The same concern appears in LLM-based code generation and
repository-level issue solving. EvalPlus strengthens code-generation
benchmarks with additional automatically constructed tests and shows
that reported \texttt{pass@k} values and model rankings can change when
the accepted solutions are subjected to a stronger test oracle
\cite{liu2023evalplus}. PatchDiff examines patches judged plausible by
SWE-bench Verified. Running a more complete collection of developer
tests reveals failures in 7.8\% of accepted patches, while generated
differential tests detect behavioral discrepancies between 29.6\% of the
examined patches and developer reference patches. Because discrepancies
can arise from legitimate implementation differences, the authors
additionally investigate correctness by manual analysis and estimate an
average inflation of 6.4 percentage points in reported issue-resolution
rates \cite{wang2026patchdiff}. The separation between \emph{behavioral
difference} and \emph{confirmed incorrectness} is important: a differential test supplies evidence of disagreement, and a verdict on correctness requires further analysis.

SWE-ABS provides another independent examination of positive judgments
in SWE-bench Verified. It strengthens tests using coverage-driven
augmentation and mutation-based adversarial construction, then
re-evaluates patches that had previously passed the benchmark. The study
reports that 19.78\% of these formerly accepted patches are rejected by
the strengthened tests; the leading agent's reported score decreases
from 78.80\% to 62.20\%, moving it from first to fifth place
\cite{yu2026sweabs}. The synthetic mutants are used to expose weaknesses
in test suites, not to measure the rate at which the original benchmark
wrongly rejected correct patches. SWE-ABS therefore supplies evidence
about accepted solutions that fail under stronger tests, while leaving
the population of previously rejected solutions outside its primary
re-evaluation.

A second group concerns rejection decisions and evaluations of C-to-Rust
translations themselves. Evidence about unsupported rejections is less
common in the studies reviewed here. In a cross-language investigation that does not include
Rust, Rabbi et al.~identify evaluation-related failure modes including
missing linker flags, incorrect compilation commands, unsuitable
inherited time limits, and mistakes in converted tests. They manually
inspect a subset of failing translations, showing that some reported
failures can originate in the evaluation setup rather than in the
translation itself \cite{rabbi2026beyond}. Such results demonstrate that
the failure side merits scrutiny, although they do not yield a
cross-system estimate of unsupported test-stage rejection decisions for
C-to-Rust translations. Furthermore, build-configuration failures are a
different population from the test-stage rejections examined in the
present thesis.

Within C-to-Rust translation, several studies look beyond the original
testing conditions. Li et al.'s user study compares 31 Rust translations
written by participants against eight C programs using fuzzing and
documents behavioral discrepancies \cite{li2025userstudy}. Reboot uses
separately prepared validation suites to evaluate its translated
interpreters, while RustAssure uses symbolic comparison to uncover bugs
not detected by existing unit tests
\cite{wang2026reboot,bai2025rustassure}. These studies strengthen the
evidence available for their own experimental subjects, but their
primary purpose is not to audit a shared set of published PASS and
test-stage FAIL decisions from multiple translation systems.

The closest precedents examine C-to-Rust outputs directly but pursue
different objectives. ACToR offers a particularly close precedent
because it also examines
translations produced by \emph{another} C-to-Rust system. In addition to
evaluating its own translation loop, ACToR applies differential testing
and refinement to C2SaferRust translations for seven standalone
programs. Before ACToR's modifications, the C2SaferRust translations
achieve a 76.3\% pass rate on the held-out validation tests used in that
experiment; after the adversarial refinement loop, the pass rate rises
to 92.9\% \cite{actor2025}. The initial 76.3\% is a validation result
obtained in ACToR's experimental setting, not necessarily the original
PASS/FAIL rate reported by C2SaferRust. Similarly, the 92.9\% pertains
to \emph{modified} translations, not a second verdict on unchanged
submissions. ACToR thus establishes a concrete precedent for applying
additional behavioral testing to published outputs of another
translator, while retaining a different purpose from a post-hoc audit:
its testing actively drives improvements to the translations. A separate
contribution by Rustine recalculates structural \texttt{unsafe} measures
for multiple C-to-Rust systems with a common measurement procedure,
illustrating cross-system consistency for a different class of metrics
\cite{dehghan2025rustine}.

The TRACTOR program also provides independent testing and evaluation.
MIT Lincoln Laboratory's first test-and-evaluation report covers six
participating development teams and two baselines on 150 Battery 01
cases, using predefined test vectors and excluding those marked as
exhibiting known UB from automated scoring
\cite{tractor-first-evaluation-report}. This is important evidence that
C-to-Rust systems can be evaluated by a common external organization. It should, however, be described as an \emph{independent primary evaluation}, which assigns its own scores to the systems. The report and the
present thesis need not involve the same released translations or runner
revision.

\begingroup
\footnotesize
\begin{xltabular}{\textwidth}{@{}>{\raggedright\arraybackslash}p{.13\textwidth}>{\raggedright\arraybackslash}p{.16\textwidth}>{\raggedright\arraybackslash}X >{\raggedright\arraybackslash}X >{\raggedright\arraybackslash}p{.18\textwidth}@{}}
\caption{Studies relevant to assessing the validity of test-based evaluation results.}\label{tab:rw-reevaluation}\\
\toprule
\textbf{Study} & \textbf{Role} & \textbf{Decisions or outputs} & \textbf{Additional assessment} & \textbf{Effect on results} \\
\midrule\endfirsthead
\toprule
\textbf{Study} & \textbf{Role} & \textbf{Decisions or outputs} & \textbf{Additional assessment} & \textbf{Effect on results} \\
\midrule\endhead
Qi et al.\ \cite{qi2015plausibility} & Post-hoc audit & Accepted repair patches & Replayed test suites and manual review & Recounts plausible and correct repairs \\
Smith et al.\ \cite{smith2015overfitting} & Method study & Patches passing repair-time tests & Held-out independent tests & Exposes overfitting \\
EvalPlus\ \cite{liu2023evalplus} & Benchmark re-evaluation & Accepted generated code & Augmented tests & Changes pass@k and ranking \\
PatchDiff\ \cite{wang2026patchdiff} & Post-hoc audit & Accepted SWE-bench patches & Developer tests, differential tests, manual review & Estimates inflated resolution rates \\
SWE-ABS\ \cite{yu2026sweabs} & Benchmark re-evaluation & Previously accepted SWE-bench patches & Coverage- and mutation-based test strengthening & Revises rates and rankings \\
Rabbi et al.\ \cite{rabbi2026beyond} & Failure analysis & Failed cross-language translations & Manual examination of build and test setup & Identifies rejection mechanisms \\
Li et al.\ \cite{li2025userstudy} & Behavioral study & Participant-written C-to-Rust programs & Differential fuzzing against C & Reports discrepancies \\
ACToR\ \cite{actor2025} & External-output assessment and refinement & Published C2SaferRust translations & Held-out tests, differential repair & 76.3\% before; 92.9\% after repair \\
MIT LL\ \cite{tractor-first-evaluation-report} & Independent primary evaluation & Multi-team C-to-Rust translations & Fixed TRACTOR vectors & Reports its own scores \\
\bottomrule
\end{xltabular}
\endgroup

These precedents establish that independent or strengthened evaluation
can materially alter the interpretation of software-generation results.
They also delineate distinct aims: rechecking a published decision,
constructing an independent initial score, and improving an output using
newly obtained tests.

\section{Research Gap and Thesis Positioning}
\label{sec:rw-positioning}

The preceding work establishes both the usefulness and the limitations
of the evidence commonly used to evaluate C-to-Rust translations.
Compiling and passing a benchmark suite are operationally meaningful
outcomes, but the original tests may fail to expose divergences or may
reject programs because of inappropriate expectations. Differential
testing, symbolic execution, and formal techniques provide additional
ways to examine behavior, but each brings its own observation
assumptions, input restrictions, or verification bounds. Research in
program repair, code generation, and issue solving further demonstrates
that the validity of published acceptance decisions can be studied
independently of the systems that produced the outputs.

The present thesis does not claim novelty for differential testing,
evaluating translations produced by another system, or recalculating
success rates in themselves. In particular, ACToR has already applied
additional behavioral tests to published C2SaferRust translations and
improved their validation performance through further repair
\cite{actor2025}. PatchDiff and SWE-ABS have independently reassessed
positive benchmark decisions and their effects on reported results in
repository-level issue solving \cite{wang2026patchdiff,yu2026sweabs}.
MIT Lincoln Laboratory has conducted a standardized independent
evaluation of multiple C-to-Rust systems
\cite{tractor-first-evaluation-report}. These are substantive
precedents, but their aims and measurement populations differ from those
of this thesis.

Within the surveyed literature, the remaining gap is a \emph{combined post-hoc assessment} of the test-based decisions published for several C-to-Rust translation systems. The studies reviewed above each cover part of it. Work on test adequacy and on differential testing shows that an accepted translation can still diverge from its C program (Sections~\ref{sec:rw-compilation} and~\ref{sec:rw-differential}), yet none of the studies reviewed here measures how often this holds among published PASS decisions; this is the subject of RQ1. Reports of unsuitable test expectations and analyses of failure causes show that a rejection can originate in the evaluation itself (Sections~\ref{sec:rw-compilation} and~\ref{sec:rw-reevaluation}), yet none estimates the share of unsupported test-stage rejections for C-to-Rust systems; this is the subject of RQ2. Studies that reassess benchmark results show that reclassified decisions change the reported rates (Section~\ref{sec:rw-reevaluation}). For C-to-Rust translation, that effect has not been computed over a shared population in which the decisions that could not be re-examined remain visible; this is the subject of RQ3.

The thesis addresses this gap on the translations and decisions released in the HARVEST results repository for the TRACTOR Public Test Corpus~\cite{harvest-results,tractor-test-corpus,schesch2026actor}. Chapter~\ref{ch:methodology} presents the design of the study.
